% ==============================================================================
% MÓDULO PRE-CAPÍTULOS: INTRODUCCIÓN GENERAL (ENFOQUE CUALITATIVO FENOMENOLÓGICO)
% Archivo: pre_caps/09_introduccion.tex
% Secuencia de Contextualización: Local (CS Belén) -> Nacional (Perú) -> Mundial (Global)
% Plantilla base: PROTOCOLO PROYECTOS CUALITATIVOS FINAL (UNSCH - EN 486)
% ==============================================================================
\section*{INTRODUCCIÓN}
\label{sec:intro}
\addcontentsline{toc}{section}{INTRODUCCIÓN}

La calidad de atención y el cuidado humano en el primer nivel de atención constituyen un imperativo ético y disciplinar indispensable para garantizar la dignidad de la persona y la consolidación de la cobertura sanitaria universal. En el marco de la investigación cualitativa bajo el diseño fenomenológico, la comprensión de este fenómeno no parte de abstracciones teóricas desvinculadas, sino de la inmersión vivencial y situada en los escenarios concretos donde acontece el encuentro interpersonal entre el profesional de salud y el usuario doliente. En estricta concordancia con la guía oficial para formular proyectos cualitativos de la Escuela Profesional de Enfermería de la Universidad Nacional de San Cristóbal de Huamanga (UNSCH), a cargo del Dr. Manglio Aguirre Andrade, la fundamentación del presente estudio se articula a través del recorrido inverso de contextualización (del ámbito específico al general) y el desarrollo riguroso de sus dimensiones justificatorias:

\paragraph{1. Contextualización de la problemática en tres niveles (Recorrido Inverso):}
\begin{itemize}
    \item \textbf{Nivel Local (Centro de Salud Belén, Ayacucho):} En el plano local y de manera estricta y exclusiva, la investigación se sitúa en el Centro de Salud Belén (Categoría I-3, Microred Huamanga), establecimiento cabecera que atiende a una población adscrita de 18,450 habitantes en el distrito de Ayacucho, a 2,760 metros sobre el nivel del mar. La realidad cotidiana del establecimiento está marcada por una profunda vulnerabilidad: el 68.2\% de los usuarios vive en condición de pobreza según el SISFOH, más del 94.5\% depende del Seguro Integral de Salud (SIS) y el 72\% posee bilingüismo activo quechua chanka-castellano. Para esta población, la calidad del cuidado no reside en la frialdad de un trámite clínico, sino en la necesidad vital de un trato digno, acogedor y respetuoso expresado en pautas andinas como el saludo afectuoso (\textit{napaykuy}), la hospitalidad en la atención (\textit{allin chaskiy}) y el respeto mutuo (\textit{respetanakuy}). El estudio responde de forma vinculante a la \textbf{Prioridad Regional 10 de Investigación en Salud de Ayacucho 2025--2030}, aprobada por la Dirección Regional de Salud (DIRESA) de Ayacucho \cite{diresa2025}: \textit{``Organización, gestión, calidad y accesibilidad de los servicios de salud''}.
    
    \item \textbf{Nivel Nacional (Perú):} En el contexto del país, la problemática local refleja las tensiones estructurales del sistema sanitario peruano, donde más del 78\% de las consultas ambulatorias recaen en el primer nivel de atención (categorías I-1 a I-4). La Política Nacional de Calidad en Salud del Ministerio de Salud (MINSA) \cite{minsacalidad2021} reconoce que la persistencia de barreras de trato interpersonal, la prisa protocolar con consultas abreviadas de escasos minutos y la debilidad en la comunicación empática en lenguas originarias erosionan la confianza ciudadana y la adherencia a los tratamientos. Esta situación motivó al Estado peruano a priorizar la investigación en servicios sanitarios mediante la \textbf{Línea 10 de Investigación en Salud al 2030}: \textit{``Sistemas y servicios de salud, acceso y cobertura universal''} (Resolución Ministerial N° 424-2025/MINSA) \cite{minsainv2025}, con el fin de generar evidencia empírica directa para humanizar la atención asistencial.
    
    \item \textbf{Nivel Mundial (Internacional o Global):} En el plano internacional, los organismos rectores como la Organización Mundial de la Salud (OMS), la OCDE y el Banco Mundial \cite{who2018} advierten que los sistemas sanitarios del mundo atraviesan una severa crisis de deshumanización asistencial provocada por la sobrecarga burocrática y el predominio de un paradigma biomédico tecnicista. Este enfoque reduce el acto clínico a procedimientos mecánicos de producción asistencial, invisibilizando la dimensión afectiva, el respeto al pudor somático y los determinantes socioculturales de la salud, lo que demanda con urgencia reorientar los servicios hacia un cuidado centrado en la persona y culturalmente seguro.
\end{itemize}

\paragraph{2. Importancia y relevancia social del proyecto:}
El estudio aborda un problema crítico de alta sensibilidad humana: la brecha latente entre los lineamientos institucionales y la vivencia sentida del usuario andino frente al cuidado de salud. Al develar los significados, frustraciones, silencios y experiencias de sosiego de las madres de familia en CRED, gestantes y personas adultas mayores del Centro de Salud Belén, se genera conocimiento de primera fuente para visibilizar la urgencia de una atención digna. Los beneficiarios directos serán los usuarios de la zona de Belén, quienes contarán con una vía científica para expresar sus demandas de respeto y buen trato; asimismo, los hallazgos beneficiarán al equipo de gestión del establecimiento, a la Red de Salud Huamanga y al colectivo profesional de enfermería.

\paragraph{3. Finalidad o valor práctico:}
A partir de sus objetivos, la investigación busca aportar evidencia empírica cualitativa para responder al problema de la despersonalización del acto de cuidado. Sus resultados permitirán proponer pautas concretas de mejora a las autoridades sanitarias competentes: el diseño de protocolos de acogida clínica humanizada con pertinencia lingüística en quechua chanka, la reorganización del tiempo de consulta para asegurar una escucha atenta y el fortalecimiento de la empatía en la relación enfermera-paciente en el Centro de Salud Belén.

\paragraph{4. Valor teórico y tipología de la investigación:}
Conforme a la tipología metodológica de la investigación científica, el estudio se clasifica como una \textbf{investigación básica o sustantiva}, con \textbf{nivel descriptivo-interpretativo} y un riguroso \textbf{diseño cualitativo fenomenológico} \cite{sampieri2018, creswell2013}. Siguiendo las directrices epistemológicas de Edmund Husserl \cite{husserl1949}, Max van Manen \cite{vanmanen1990} y Roberto Hernández-Sampieri \cite{sampieri2018}, el proyecto llena un vacío crítico de conocimiento al contrastar la base teórica existente con la vivencia del poblador andino. Para ello, articula dialógicamente el modelo del \textbf{Proceso Interpersonal de Calidad en Salud de Avedis Donabedian} \cite{donabedian1980, donabedian2005} y la \textbf{Teoría del Cuidado Humano de Jean Watson} \cite{watson2008} con los cuatro existenciales universales del \textit{Lebenswelt}: relacionalidad (el encuentro humano y diálogo en quechua), temporalidad (tiempo y calma de la consulta), espacialidad (privacidad del consultorio) y corporalidad (respeto al pudor físico y alivio del dolor \textit{nanay}).

\paragraph{5. Valor metodológico:}
En el ámbito metodológico, el proyecto aporta una guía de entrevista fenomenológica en profundidad semiestructurada bilingüe y una matriz de categorización temática cualitativa, centradas exclusivamente en la calidad del cuidado. Estos instrumentos y procedimientos se ajustan a los cuatro criterios de rigor cualitativo establecidos por Guba y Lincoln \cite{gubalincoln1985}: credibilidad, transferibilidad, consistencia y confirmabilidad, complementados con la técnica de triangulación múltiple de métodos y fuentes de información \cite{okuda2005}.

\paragraph{6. Justificación institucional y bioética:}
El protocolo cumple los estándares científicos de la cátedra de EN 486 (fundamentados en el rigor epistemológico de Mario Bunge \cite{bunge2014}), respetando de forma irrestricta las normas bioéticas internacionales de la Declaración de Helsinki \cite{helsinki2013} y las pautas éticas del CIOMS \cite{cioms2016}, garantizando el consentimiento informado voluntario, la confidencialidad, el anonimato mediante códigos alfanuméricos y la no maleficencia para con los participantes.

\newpage
